\section{How desks review forecast models}

Our backtest pack (\texttt{reports/figures/backtests/}) copies the review
workflow of a real trading desk. This note explains that workflow and how
to read each chart. Interview-critical: knowing WHAT to check signals
production experience better than any model architecture talk.

\subsection{The three review loops}

\begin{enumerate}
  \item \textbf{Daily (automated).} Every morning: yesterday's forecast
        vs what actually happened, vs the benchmark. One number, one
        chart, no meeting. Our daily cron does this for load and now for
        price (\texttt{reports/daily/}). The realized price lands the next
        day and is plotted over the published band.
  \item \textbf{Weekly/monthly (the model review).} The six-chart pack.
        Questions, in order: is it drifting? is the edge real? where does
        it fail? can risk trust the band? what did the worst days look
        like? is it biased?
  \item \textbf{Promotion gate.} A challenger runs in shadow mode for a
        fixed window (ours: 14 days). It is scored daily but its numbers
        are not official. It replaces the incumbent only if it wins over
        the whole window on metrics agreed BEFORE the trial. This kills
        cherry-picking. Logged in \texttt{DECISIONS.md}.
\end{enumerate}

\subsection{Reading the six charts}

\begin{itemize}
  \item \textbf{Rolling MAE (drift monitor).} A model trained on last
        year quietly rots when the market changes (new interconnector,
        new capacity, fuel regime). Healthy: stable gap below the naive.
        Sick: the gap closes. Ours holds $\sim$10 EUR/MWh of gap for two
        years.
  \item \textbf{Cumulative edge.} Sum of daily (naive MAE $-$ model MAE)
        over time --- read like a P\&L curve. A steady climb means the
        edge is earned every week. A staircase means one lucky event.
        Desks distrust anything they cannot see accrue.
  \item \textbf{Error by delivery hour.} Prices have physics per hour:
        night is flat, midday is solar, hour 19--20 is the evening ramp
        where scarcity bites. Everyone's MAE triples at hour 19. If a
        model is flat across hours, suspect it is ignoring the market
        structure.
  \item \textbf{Quantile calibration.} If the P90 is pierced 26\% of the
        time it is not a P90, and any position sized on it is underhedged.
        Risk teams read this chart before MAE. Our LightGBM P50 wins MAE
        but its band fails here --- so LEAR publishes the daily band and
        LightGBM waits for conformal calibration. That decision IS the
        production discipline.
  \item \textbf{Worst-day post-mortems.} Averages hide the days that cost
        money. Desks re-play every worst day: what did the model see, what
        happened? Ours: three winter evening spikes above every band plus
        one negative-price solar day. Common cause: no scarcity signals
        (outages, fuel) in the features yet. An explainable miss with a
        feature-roadmap answer beats an unexplained small error.
  \item \textbf{Monthly bias.} Signed error. MAE-fine but always-negative
        means systematically underpricing --- a desk position leaks money
        one way. Ours: $-2$ to $-4$ EUR/MWh, worst in winter. Candidate
        for a bias term after fuel features land.
\end{itemize}

\subsection{Metric conventions for prices}

\begin{itemize}
  \item \textbf{No MAPE.} 798 negative-price hours in our sample;
        division near zero is garbage.
  \item \textbf{rMAE} $=$ MAE / MAE(naive-1d). Below 1.0 = skill. The EPF
        literature's standard (Lago et al.\ 2021). Ours: 0.638 (LightGBM),
        0.660 (LEAR).
  \item \textbf{Pinball loss} per quantile, \textbf{coverage} for the
        band, \textbf{spike metrics} on the top-5\% hours separately.
\end{itemize}

Interview line: ``MAE picks the champion, calibration picks what you
publish, post-mortems pick the next feature. Three different charts,
three different decisions.''
